\section{六自由度刚体动力学}
\label{sec:dynamics}

整机作为刚体在三维空间中的运动由质心平动和绕质心的转动两部分组成。
本节建立完整的动力学和运动学方程，并讨论数值积分策略。

本章沿用飞行动力学文献常见的机体符号$\{\mathcal B\}$，在数值实现中它与
第\ref{sec:mapping}节的推进模型系$\{P\}$一致；固件任务系$\{F\}$须经
式(\ref{eq:frame_permutation})置换后才能与本章分量逐项比较。

\subsection{牛顿-欧拉方程}

将飞行器视为刚体（弹性变形在概念设计阶段可忽略），其六自由度动力学由
牛顿第二定律（平动）和欧拉方程（转动）在机体系中联合描述\cite{stevens2015aircraft}。

\textbf{平动方程}。质心速度$\bm{v} = [u, v, w]^T$定义在机体系中。
然而牛顿第二定律$m\dot{\bm{v}}_{\mathcal{I}} = \sum \bm{F}$中的加速度是惯性系中的绝对加速度。
利用动系求导的坐标转换公式，机体系中观测的速度变化率加上因坐标系旋转带来的牵连项
等于惯性系绝对加速度在机体系的分量：
\begin{equation}
    m \left( \dot{\bm{v}} + \bm{\omega} \times \bm{v} \right)
    = \bm{F}_{\text{推力}} + \bm{F}_{\text{气动}} + m \bm{g}_b
    \label{eq:newton_body}
\end{equation}
其中$\bm{\omega} = [p, q, r]^T$为机体角速度。$\bm{\omega} \times \bm{v}$项
反映了即使$\bm{v}$在机体系中不变，因机体旋转导致$\bm{v}$在惯性系中的方向改变
所需的向心加速度。这是机体系中牛顿方程特有的附加项，在惯性系中写牛顿方程则不需要。

\textbf{转动方程}。对机体系原点（质心）取矩，并考虑旋转部件（电机+螺旋桨）
的角动量效应：
\begin{equation}
    \bm{I} \dot{\bm{\omega}} + \bm{\omega} \times (\bm{I} \bm{\omega})
    + \bm{\omega} \times \bm{h}_{\text{转子}}
    = \bm{M}_{\text{推力}} + \bm{M}_{\text{气动}}
    \label{eq:euler_body}
\end{equation}
各物理项的含义依次为：
\begin{itemize}
    \item $\bm{I} \dot{\bm{\omega}}$：机体角加速度引起的惯性力矩（转动惯量乘角加速度）。
          对于主轴坐标系，$\bm{I} = \operatorname{diag}(I_x, I_y, I_z)$；
    \item $\bm{\omega} \times (\bm{I} \bm{\omega})$：因机体角速度矢量与角动量矢量
          不共线（除非绕主轴旋转）而产生的陀螺进动力矩。这一项引起三轴之间的动态耦合——
          例如滚转（$p$）和偏航（$r$）均不为零时，机身绕$y$轴产生$(I_x-I_z)pr$的惯性耦合力矩；
    \item $\bm{\omega} \times \bm{h}_{\text{转子}}$：转子总角动量
          $\bm{h}_{\text{转子}} = \sum_i J_{p,i} \omega_i \bm{n}_i$
          因机体旋转产生的陀螺力矩（已在第\ref{sec:propulsion}节详细讨论）。
          注意此处是$+\bm{\omega} \times \bm{h}$（式\ref{eq:euler_body}左侧），
          与式(\ref{eq:gyro_moment})的$-\bm{\Omega} \times \bm{h}$看似符号相反，
          实则为同一物理量在方程两侧的不同表述——左侧是因，右侧是果。
\end{itemize}

\textbf{重力在机体系的分量}。重力加速度在惯性系中恒为$[0, 0, g]^T$（NED约定，$z_{\mathcal{I}}$向下）。
通过姿态旋转矩阵的逆变换（从惯性系到机体系）得：
\begin{equation}
    \bm{g}_b = \bm{R}_{\mathcal{I}\mathcal{B}}
    \begin{bmatrix} 0 \\ 0 \\ g \end{bmatrix}
    = \bm{R}_{\mathcal{B}\mathcal{I}}^T
    \begin{bmatrix} 0 \\ 0 \\ g \end{bmatrix}
    \label{eq:gravity_body}
\end{equation}

\textbf{气动力建模}。在概念级仿真中，气动力可用简化的集总参数模型\cite{beard2012small}：
动压$q_\infty = \frac{1}{2}\rho V^2$（$V$为空速），迎角$\alpha = \arctan(w/u)$，
侧滑角$\beta = \arcsin(v/V)$。气动力的完整建模需要气动导数数据库（通常来自CFD或风洞试验），
不在纯推力矢量分析的范围内。在推力矢量构型的初步分析中，可先忽略气动力以聚焦推力矢量的
控制特性，随后逐步引入气动项以评估气动-推力矢量的交互效应。

\subsection{姿态运动学：四元数表示}

采用四元数$\bm{q} = [q_0, q_1, q_2, q_3]^T$（标量在前约定）
描述机体到惯性系的旋转。选择四元数而非欧拉角的主要理由是在$\theta = \pm 90^\circ$处
欧拉角存在万向锁奇异——3-2-1欧拉角序列在$\theta = \pm 90^\circ$时，
$\phi$和$\psi$描述同一物理旋转自由度，导致运动学方程退化。
四元数无此奇异点，对于可能包含大俯仰角的推力矢量构型尤其必要\cite{markley2014fundamentals,kuipers1999quaternions}。
论文记号与JavaScript对象字段的对应关系为$q_0=q_w$、$q_1=q_x$、$q_2=q_y$、$q_3=q_z$；
代码以\texttt{\{x,y,z,w\}}存储，但四元数乘法和旋转语义与本章标量在前的公式一致。

姿态运动学方程（机体系角速度$\bm{\omega}_b = [p, q, r]^T$驱动四元数演化）：
\begin{equation}
    \dot{\bm{q}} = \frac{1}{2} \bm{q} \otimes
    \begin{bmatrix} 0 \\ p \\ q \\ r \end{bmatrix}
    \label{eq:quaternion_kinematics}
\end{equation}
写成矩阵形式（便于数值实现）：
\begin{equation}
    \begin{bmatrix} \dot{q}_0 \\ \dot{q}_1 \\ \dot{q}_2 \\ \dot{q}_3 \end{bmatrix}
    = \frac{1}{2}
    \begin{bmatrix}
        0 & -p & -q & -r \\
        p &  0 &  r & -q \\
        q & -r &  0 &  p \\
        r &  q & -p &  0
    \end{bmatrix}
    \begin{bmatrix} q_0 \\ q_1 \\ q_2 \\ q_3 \end{bmatrix}
    \label{eq:quaternion_matrix}
\end{equation}

四元数到方向余弦矩阵（$\{\mathcal{B}\}\to\{\mathcal{I}\}$）的转换：
\begin{equation}
    \bm{R}_{\mathcal{B}\mathcal{I}} =
    \begin{bmatrix}
        1-2(q_2^2+q_3^2) & 2(q_1 q_2 - q_0 q_3) & 2(q_1 q_3 + q_0 q_2) \\
        2(q_1 q_2 + q_0 q_3) & 1-2(q_1^2+q_3^2) & 2(q_2 q_3 - q_0 q_1) \\
        2(q_1 q_3 - q_0 q_2) & 2(q_2 q_3 + q_0 q_1) & 1-2(q_1^2+q_2^2)
    \end{bmatrix}
    \label{eq:dcm_from_quaternion}
\end{equation}

Web实现为保持既有显示、SAS整定和回归基线，采用下式从DCM提取兼容显示角：
\begin{equation}
    \begin{aligned}
        \phi &= \operatorname{atan2}(R_{23}, R_{33}) \\
        \theta &= -\arcsin(R_{13}) \\
        \psi &= \operatorname{atan2}(R_{12}, R_{11})
    \end{aligned}
    \label{eq:euler_from_dcm}
\end{equation}

式(\ref{eq:euler_from_dcm})不是由$\bm R_{\mathcal{B}\mathcal{I}}$提取标准
3-2-1航空欧拉角的通式。特别地，绕$+y_b$的纯正向旋转在当前实现中得到
$\theta<0$；该$\theta=-\arcsin(R_{13})$约定属于行为保持红线。论文在讨论
Web历史SAS和HUD曲线时沿用这一显示量，而在飞控任务轴、物理姿态和控制律推导中不应
把它与标准航空俯仰角混为一谈。四元数状态本身仍表示机体到惯性系的旋转，
式(\ref{eq:dcm_from_quaternion})及其向量旋转语义不因这一显示约定而改变。

\textbf{四元数的数值维护}：每次积分步后，四元数因浮点舍入误差的积累会逐渐偏离单位模长
（$\|\bm{q}\| \neq 1$）。若不加处理，范数漂移会导致DCM失去正交性，
进而使坐标变换产生系统性畸变。解决方案是对每一积分子步后的四元数执行归一化：
$\bm{q} \leftarrow \bm{q} / \|\bm{q}\|$。这是飞行仿真中的标准做法。

\subsection{惯性系位置与速度}

机体系速度经姿态变换到惯性系后积分得到惯性系位置：
\begin{equation}
    \dot{\bm{p}}_{\mathcal{I}} = \bm{R}_{\mathcal{B}\mathcal{I}} \bm{v}_{\mathcal{B}}
    \label{eq:position_integration}
\end{equation}

\subsection{数值积分方法与子步调度}

\subsubsection{积分方法选择}

对刚体六自由度系统，状态向量$\mathbf{x}=[\mathbf{v}_{\mathcal{B}}^T,\boldsymbol{\omega}_{\mathcal{B}}^T,\mathbf{q}^T]^T$（共10维）的演化由以下常微分方程描述：
\begin{align}
    \dot{\mathbf{v}}_{\mathcal{B}} &= \mathbf{F}_{\text{total}}/m - \boldsymbol{\omega}\times\mathbf{v}_{\mathcal{B}} \\
    \dot{\boldsymbol{\omega}}_{\mathcal{B}} &= \mathbf{I}^{-1}\left[\mathbf{M}_{\text{total}} - \boldsymbol{\omega}\times(\mathbf{I}\boldsymbol{\omega}) - \boldsymbol{\omega}\times\mathbf{h}_{\text{rotor}}\right] \\
    \dot{\mathbf{q}} &= \frac{1}{2}\,\mathbf{q}\otimes[0,\boldsymbol{\omega}_{\mathcal{B}}]^T
\end{align}

采用\textbf{经典四阶显式Runge-Kutta法（RK4）}联合积分全部状态，每子步执行一次完整的四阶递推：
\begin{equation}
    \begin{aligned}
        \mathbf{k}_1 &= h\,\mathbf{f}(\mathbf{x}_k) \\
        \mathbf{k}_2 &= h\,\mathbf{f}(\mathbf{x}_k + \mathbf{k}_1/2) \\
        \mathbf{k}_3 &= h\,\mathbf{f}(\mathbf{x}_k + \mathbf{k}_2/2) \\
        \mathbf{k}_4 &= h\,\mathbf{f}(\mathbf{x}_k + \mathbf{k}_3) \\
        \mathbf{x}_{k+1} &= \mathbf{x}_k + (\mathbf{k}_1 + 2\mathbf{k}_2 + 2\mathbf{k}_3 + \mathbf{k}_4)/6
    \end{aligned}
    \label{eq:rk4}
\end{equation}
其中$\mathbf{f}$为上述ODE的右侧联合函数，在RK4各阶段内由中间状态计算
$\dot{\mathbf{v}}$、$\dot{\boldsymbol{\omega}}$和$\dot{\mathbf{q}}$，同时重新计算
重力投影与气动力。推进力/力矩和转子角动量在单个物理子步内采用零阶保持；
气动力则由每个RK4阶段的中间速度与角速度重新评估。

\textbf{选择RK4的理由}：
\begin{enumerate}[label=(\arabic*)]
    \item \textbf{精度跃升}：显式欧拉为一阶精度，RK4为四阶精度（局部$O(h^5)$）。
          当前可重复回归工况中，RK4误差为$1.4979\times10^{-12}$，显式欧拉误差为
          $3.99999\times10^{-5}$；该指标用于比较离散误差，不等同于跨工况的姿态角漂移；
    \item \textbf{稳定性域更大}：显式欧拉的绝对稳定条件为$|1+h\lambda|<1$（$\lambda<0$时$h<2/|\lambda|$），RK4的稳定边界约$h<2.78/|\lambda|$；
    \item \textbf{四元数流形保护}：四元数定义在$S^3$单位球面上，RK4每阶段均执行保范投影抑制高阶误差累积；
    \item \textbf{角加速度估计兼容性}：增量非线性动态逆（第\ref{sec:INDI}节）依赖
          $\dot{\boldsymbol{\omega}}_0$估计，高精度状态传播可降低数值误差；实际闭环质量仍受
          采样、滤波群时延、传感器噪声与算法实现一致性共同限制。
\end{enumerate}

\subsubsection{子步调度}

主循环以屏幕刷新率（通常16--17\,ms）或固定频率运行，物理积分采用更小的子步。设主循环步长为$\Delta t$，取$n = \lceil \Delta t / \Delta t_{\max} \rceil$个子步（$\Delta t_{\max}=4$\,ms），每子步以$h = \Delta t / n$推进一个完整的RK4步。

子步内的计算流程为：
\begin{enumerate}[label=(\arabic*)]
    \item 从四元数提取显示欧拉角，计算SAS修正（更新积分状态）；
    \item 推进系统与六维映射（更新电机状态）；
    \item 计算子步起点气动力，随后在RK4各中间阶段按阶段状态重新计算；
    \item \textbf{RK4联合积分}：推进$\mathbf{v}_{\mathcal{B}}$、
          $\boldsymbol{\omega}_{\mathcal{B}}$和$\mathbf{q}$，每阶段重新计算重力与气动力，
          并执行四元数保范投影；推进力矩与转子角动量在该子步内保持；
    \item 用更新后的速度和姿态计算惯性系速度，再以一阶位置更新施加交互仿真的运动学约束。
\end{enumerate}

因此，10维RK4状态不含位置，不能称为位置也参与的完整六自由度联合RK4。
主循环参数\texttt{frameCap}=0.05\,s，物理子步参数\texttt{maxStep}=0.004\,s；
二者均来自\texttt{models/aircraft-model.json}的MODEL-DEFAULT设置，未经试验标定。
